\section{The actual final bootstrap in Weil II}
\label{weil-ii-dyadic-bootstrap}

The proof of the curve theorem in \emph{Weil II}, 3.2.3--3.2.15,
uses a more specific mechanism than tensor amplification alone. A local
calculation makes the weights of every relevant nonconstant constituent
integral. An earlier analytic theorem excludes weight \(2\) and all larger
weights. Consequently these integral weights are at most \(1\). The
remaining constituents are geometrically constant on the projective
line, where their first cohomology vanishes. These two exclusions make a
product-square argument improve a universal error \(e\) to \(e/2\).
The proof below retains the schemes, coefficient sheaves, boundary
terms, maps, twists, and indices which make that improvement possible.

The primary source is Pierre Deligne,
\href{https://www.numdam.org/item/PMIHES_1980__52__137_0/}
{\emph{La conjecture de Weil II}}, Publications math\'ematiques de
l'IH\'ES 52 (1980), 137--252: the introduction, printed pp.~140--141;
1.8.4, pp.~175--176; 2.2.10, p.~196; 3.1, pp.~197--200;
3.2, pp.~200--204; and 3.3, pp.~204--207. All numbered results cited
in this section are Deligne's published results unless a displayed tag
begins with WDB. The local current S20 English and French texts were
read alongside the controlling published French scan. This section is
an exposition of that proof and its exact algebraic steps, not a claim
that a construction over the absolute base has already acquired the
finite-field sheaves in this proof.

\subsection{Conventions, original weight, and the universal assertion}

Fix a finite field \(\mathbb F_q\), an algebraic closure
\(\overline{\mathbb F}_q\), a prime \(\ell\ne\operatorname{char}
\mathbb F_q\), and Deligne's coefficient isomorphism
\(\iota:\overline{\mathbb Q}_\ell\longrightarrow\mathbb C\).
For a scheme or sheaf with subscript \(0\), omission of that subscript
means extension to \(\overline{\mathbb F}_q\). Use geometric
Frobenius and the weight convention
\[
 w_q(\alpha)=2\log_q|\iota\alpha|,
 \qquad w_{q^d}(\alpha^d)=w_q(\alpha),
 \qquad w_q(\alpha^r)=r w_q(\alpha).
 \tag{WDB1}
\]
A Tate twist \((a)\) multiplies Frobenius eigenvalues by
\(q^{-a}\) and lowers their weights by \(2a\).

Let \(U_0\) be a smooth curve and let \(\mathcal F_0\) be a lisse
pointwise \(\iota\)-pure sheaf of weight \(\beta\) on \(U_0\).
Deligne proves the universal assertions
\[
 \mathsf B_k:\quad
 w_q(\alpha)\le \beta+1+2^{-k}
 \quad\text{for every eigenvalue of }F
 \text{ on }H_c^1(U,\mathcal F),
 \qquad k=0,1,2,\ldots .
 \tag{WDB2}
\]
The quantifiers range over all such curves, sheaves, weights, and finite
fields. This universality is used when the induction is applied to new
coefficient sheaves on the base of a pencil, rather than only to the
original \(\mathcal F_0\).

Here is the precise relation to the weight-zero formulation of 3.2.4.
Choose \(c\in\overline{\mathbb Q}_\ell^*\) with
\(\iota(c)=q^{-\beta/2}\), and let \(\mathcal L_c\) be the
constant geometric Weil sheaf whose arithmetic datum is Frobenius
\(c\). Deligne's twisting operation 1.2.7 gives
\[
 \mathcal F_0^{[0]}=\mathcal F_0\otimes\mathcal L_c,
 \quad
 H_c^1(U,\mathcal F^{[0]})
   =H_c^1(U,\mathcal F)\otimes\mathcal L_c,
 \quad
 \alpha^{[0]}=c\alpha,
 \quad w_q(\alpha^{[0]})=w_q(\alpha)-\beta.
 \tag{WDB3}
\]
Thus the original weight and eigenvalue are recovered by tensoring with
\(\mathcal L_{c^{-1}}\) and multiplying by \(c^{-1}\). We use
\(\mathcal F_0^{[0]}\), not a change in the meaning of
\(\mathcal F_0\), for the auxiliary weight-zero proof. This twist is
made in the Weil-sheaf setting allowed by Deligne; it is not an assertion
that every such character extends to an \(\ell\)-adically continuous
character of the profinite arithmetic Galois group.

We will repeatedly use the following elementary linear algebra. If a
linear automorphism preserves a finite filtration, a basis adapted to
the filtration makes its matrix block triangular. Its characteristic
polynomial is the product of the characteristic polynomials on the
successive quotients. Hence a uniform upper weight bound is inherited
by subspaces, quotients, and finite extensions. For a convergent
Frobenius-equivariant spectral sequence, each term on a later page is
a subquotient of an earlier term; the abutment has a filtration with
the \(E_\infty\) terms as successive quotients. These facts prove the
weight comparisons below without assuming spectral-sequence
degeneration or Frobenius semisimplicity.

\subsection{The two earlier inputs: a strict bound and local integral shifts}

Deligne's 1.8.2 supplies the initial estimate
\(w_q(\alpha)\le\beta+2\), hence \(\mathsf B_0\).
For a smooth affine curve it applies directly with the open immersion
in 1.8.2 taken to be the identity: the Euler product and the
\(H_c^2\) denominator have no zero in
\(|t|<q^{-(\beta+2)/2}\), so the numerator has none there.
For a projective component, remove a closed point and use the
surjection from the compactly supported first cohomology of the
remaining affine curve, proved explicitly in (WDB7)--(WDB8) below.
Connected components are handled by a finite extension and (WDB1).

The stronger earlier result used in the induction is 2.2.10:
\[
 \alpha\in\operatorname{Spec}
   \bigl(F;H_c^1(U,\mathcal F)\bigr)
 \quad\Longrightarrow\quad w_q(\alpha)<\beta+2.
 \tag{WDB4}
\]
The inequality is strict. Deligne proves it in Section 2 using the
Hadamard--de la Vall\'ee-Poussin method and the cohomological expression
for the relevant \(L\)-function. After weight-zero twisting and
reduction to an irreducible representation, 2.2.9 excludes zeros on
the boundary line \(\operatorname{Re}s=1\); 2.2.8 and 2.2.9.1
identify this with the absence of the corresponding numerator roots
on \(|t|=q^{-1}\). The possible geometrically constant character is
accounted for by the simple pole stated there, rather than discarded.
Its pole order equals the order of the \(H_c^2\) denominator zero,
so a numerator zero at that point would contradict the exact simple
pole and is excluded as well.
Since a reciprocal root \(t=\iota(\alpha)^{-1}\) has modulus
\(q^{-w_q(\alpha)/2}\), the absence of numerator roots for
\(|t|\le q^{-1}\) is exactly \(w_q(\alpha)<2\) in weight zero.
This is an earlier proved input in the order of the paper; it is not
deduced by assuming the conclusion of 3.2.3.

The second earlier result is local monodromy purity, 1.8.4. If a lisse
sheaf on a punctured smooth curve is pointwise \(\iota\)-pure of
weight \(b\), and \(M\) is the centered local monodromy filtration
on its geometric generic stalk \(V\), then
\[
 \operatorname{Gr}_j^M V\text{ is }\iota\text{-pure of weight }b+j,
 \qquad j\in\mathbb Z.
 \tag{WDB5}
\]
In particular every Frobenius eigenvalue of the local Weil
representation has weight in \(b+\mathbb Z\). Under unipotent
inertia the graded pieces have trivial inertia: writing
\(\rho(\sigma)=\exp(t_\ell(\sigma)N)\), the operator \(N\)
lowers \(M\) by two, so its exponential acts as the identity on
every \(\operatorname{Gr}_j^M\).

The weight shift in (WDB5) is itself proved in 1.8.4, not read off
from the existence of a nilpotent operator. To recall its numerical
content, twist to weight zero and use a finite covering so that all
inertia is unipotent. Then \(\ker N=V^I\), and 1.8.1 says its
eigenvalues have modulus at most \(1\). For a primitive bottom
piece \(P_{-j}V\), Deligne's tensor decomposition 1.6.14.4 makes
\[
 P_{-j}V\otimes P_{-j}V(-j)
 \quad\text{a direct summand of}\quad P_0(V\otimes V).
 \tag{WDB6}
\]
If \(\alpha\) occurs in \(P_{-j}V\), the resulting eigenvalue
\(\alpha^2q^j\) therefore has modulus at most \(1\), giving
\(|\iota\alpha|\le q^{-j/2}\). Applying the same argument to
the dual primitive piece, whose corresponding eigenvalue is
\(\alpha^{-1}q^{-j}\) by 1.6.14.5, gives
\(|\iota\alpha|\ge q^{-j/2}\). Equality follows. The primitive
decomposition 1.6.14.3 and the Tate shifts on its strings give (WDB5).
Thus the integral shifts later used by the pencil argument have
already been tied to the actual local Weil action.

\subsection{The exact open-extension map used by the induction}

Let \(j:W_0\hookrightarrow C_0\) be a dense open immersion of
curves, with finite complement, and let \(\mathcal K_0\) be a
constructible sheaf on \(C_0\). Extension by zero gives the short
exact sequence
\[
 0\longrightarrow j_!j^*\mathcal K_0
 \longrightarrow\mathcal K_0
 \longrightarrow \mathcal K_0/j_!j^*\mathcal K_0
 \longrightarrow0.
 \tag{WDB7}
\]
The quotient is supported on finitely many points. Over the algebraic
closure its positive-degree compactly supported cohomology vanishes:
each geometric point has only degree-zero cohomology and a finite
disjoint union has the direct sum of those cohomologies. The long
exact sequence of (WDB7) consequently gives the Frobenius-equivariant
surjection
\[
 H_c^1(W,j^*\mathcal K)
 \twoheadrightarrow H_c^1(C,\mathcal K).
 \tag{WDB8}
\]
If \(W_0\) is smooth and \(j^*\mathcal K_0\) is lisse and
pointwise pure of weight \(b\), assertion \(\mathsf B_k\)
applied to the source proves that every eigenvalue of the target has
weight at most \(b+1+2^{-k}\). This is precisely 3.2.5\((k)\).
No purity of the possibly singular boundary stalks of
\(\mathcal K_0\) is assumed in this step.

\subsection{The product and pencil, with all three reduction hypotheses}

Assume \(\mathsf B_k\). Initially impose the three hypotheses of
3.2.6 on the auxiliary sheaf \(\mathcal F_0^{[0]}\).

\begin{enumerate}
\item[(A)] If \(X_0\) is the smooth projective compactification of
\(U_0\), the local monodromy of \(\mathcal F^{[0]}\) at
\(X-U\) is tame and unipotent.
\item[(B)] Put \(S_0=X_0\times X_0\),
\(V_0=U_0\times U_0\), and \(D_0=S_0-V_0\). Choose a projective
embedding and a hyperplane pencil with axis \(A_0\) satisfying
3.1.1(A)--(D): the axis is transverse to \(S\) and disjoint from
\(D\); the only singularities of the pencil on \(S\) are ordinary
quadratic points outside \(D\); its only singularities on the
normalization of \(D\) are ordinary quadratic points away from the
preimages of the nodes of \(D\); and each exceptional fibre contains
exactly one exceptional point.
\item[(C)] These exceptional points are defined over \(\mathbb F_q\).
\end{enumerate}

Write \(P_0=A_0^*\simeq\mathbb P^1_{\mathbb F_q}\) for the
parameter line of the pencil, and use the incidence surface
\[
 \widetilde S_0=\{(s,t)\in S_0\times P_0:s\in H_t\},
 \quad \pi:\widetilde S_0\longrightarrow S_0,
 \quad\bar f:\widetilde S_0\longrightarrow P_0.
 \tag{WDB9}
\]
It is the blow-up at the finite set \(S_0\cap A_0\). Put
\(\widetilde V_0=\pi^{-1}(V_0)\), let
\(\widetilde j:\widetilde V_0\hookrightarrow\widetilde S_0\)
be the open immersion, and set \(f=\bar f\widetilde j\).
The actual coefficient sheaf is
\[
 \mathcal G_0
 =\operatorname{pr}_1^*\mathcal F_0^{[0]}
   \otimes\operatorname{pr}_2^*\mathcal F_0^{[0]}
 \quad\text{on }V_0.
 \tag{WDB10}
\]
Both tensor factors are retained. In terms of the original sheaf this
is \((\mathcal F_0\boxtimes\mathcal F_0)\otimes\mathcal L_{c^2}\),
so its original weight \(2\beta\) has been recorded by an explicit
invertible twist. Its auxiliary weight is zero.

Let \(T_0\subset P_0\) be the finite set of exceptional values,
\(W_0=P_0-T_0\), and \(w:W_0\hookrightarrow P_0\). Define
\[
 \mathcal R^b_0=R^b f_!\pi^*\mathcal G_0
   =R^b\bar f_*\widetilde j_!\pi^*\mathcal G_0.
 \tag{WDB11}
\]
These sheaves are lisse over \(W_0\). The compact-support Leray
spectral sequence is Frobenius-equivariant:
\[
 E_2^{ab}=H^a(P,\mathcal R^b)
 \ \Longrightarrow\ H_c^{a+b}(\widetilde V,\pi^*\mathcal G).
 \tag{WDB12}
\]
Since both the base and the fibres have dimension one, total degree
two has only the terms \((a,b)=(0,2),(1,1),(2,0)\).

\subsection{Why the middle coefficient sheaf has pure constituents}

Let \(x\in|W_0|\). Denote the proper fibre of \(\bar f\) by
\(\widetilde S_x\), its open part by \(Y_x\), and the open
immersion by \(j_x:Y_x\hookrightarrow\widetilde S_x\). Proper
base change applied to (WDB11) gives the exact stalk identification
\[
 (\mathcal R^b_0)_{\bar x}
 \simeq H^b(\widetilde S_{\bar x},
                j_{x!}\pi^*\mathcal G)
 \simeq H_c^b(Y_{\bar x},\pi^*\mathcal G).
 \tag{WDB13}
\]
These are identifications of representations of geometric Frobenius
over \(\kappa(x)\). The open fibre is a smooth curve for
\(x\in W_0\); it need not be connected.

There is a notation defect in the published display 3.2.7.1, printed
p.~202: it writes ordinary \(H^1(Y,\pi^*\mathcal H)\) for an
open fibre while the left side is \(R^1f_!\). The controlling French
scan and both current S20 transcriptions show the same display.
Equation (WDB13) states the corrected support convention and its
base-change proof. Ordinary cohomology of the open fibre is not used
in place of compactly supported cohomology.

Set \(\mathcal H_0=\mathcal G_0\oplus\mathcal G_0^\vee\).
At a closed point the eigenvalues of \(\mathcal G_0\) have
\(\iota\)-modulus one. The images of the eigenvalues of its dual
are their inverses, hence their complex conjugates. Therefore the
characteristic polynomial of the direct sum has real coefficients,
and \(\mathcal H_0\) is pointwise pure of weight zero and
\(\iota\)-real. Deligne's 3.2.1--3.2.2 applied to the smooth fibre
in (WDB13) shows that
\(w^*R^1f_!\pi^*\mathcal H_0\) is \(\iota\)-real.
The sheaf \(w^*\mathcal R^1_0\) is an actual direct summand of
it because pullback and \(R^1f_!\) preserve finite direct sums.
The purity criterion 1.5.1 now gives a finite filtration of
\(w^*\mathcal R^1_0\) whose successive quotients are lisse and
pointwise \(\iota\)-pure. This proves 3.2.7 with its coefficient
and support maps specified.

To make this use of 3.2.1 explicit, let \(\mathcal E_0\) be
pointwise \(\iota\)-pure of weight \(b\) and \(\iota\)-real on
a smooth geometrically connected curve \(C_0\). Put
\(P_a(t)=\iota\det(1-Ft,H_c^a(C,\mathcal E))\).
Grothendieck's formula says that \(P_1/(P_0P_2)\) has real Taylor
coefficients and hence is a real rational function. The roots of
\(P_2\) all have modulus \(q^{-(b+2)/2}\) by 1.4.3; those of
\(P_0\), if any, have modulus \(q^{-b/2}\); and (WDB4) excludes
roots of \(P_1\) on the former circle. Consequently the poles on
that circle, with their multiplicities, are exactly the roots of
\(P_2\). They are stable under complex conjugation, so \(P_2\)
has real coefficients and constant term one. If \(C\) is affine,
\(P_0=1\). If it is projective, the correct scalar duality is
\[
 H^0(C,\mathcal E)^\vee
    \simeq H_c^2(C,\mathcal E^\vee(1)).
 \tag{WDB13a}
\]
The coefficient sheaf on the right is again \(\iota\)-real and
pure, so the just-proved degree-two argument and inversion of
eigenvalues show that \(P_0\) is real. It follows that \(P_1\)
is real. If geometric components are permuted in a cycle of length
\(d\), their cohomological characteristic polynomial is
\(\det(1-F^d t^d,H_c^a(C',\mathcal E))\) for one component
\(C'\); the preceding argument over \(\mathbb F_{q^d}\)
makes this polynomial real, proving the extension 3.2.2.

There is a second printed notation defect in the proof of 3.2.1,
p.~200: it writes \(\mathcal E^\vee(-1)\) in place of the
\(\mathcal E^\vee(1)\) in (WDB13a). The printed group equals
\(H^0(C,\mathcal E)^\vee(-2)\), so it still implies the required
reality after the real scalar factor \(q^2\) is retained. Formula
(WDB13a) records the exact untwisted duality used here; the later
3.2.15 itself has the correct \((1)\) twist.

\subsection{Vanishing cycles detect the nonconstant weights}

The calculation in 3.1 treats exactly three exceptional geometries:
an ordinary node of the fibre away from \(D\), tangency of the fibre
to one smooth branch of \(D\), and passage of the fibre through a
node of \(D\). Let \(\bar t\in T\) and \(\bar x\) be its
unique exceptional point. The vanishing-cycle sheaves are supported
at \(\bar x\) and vanish in every degree other than one. If
\(\varepsilon(B)\) is the one-dimensional sign representation
of the two branches or two local sheets, the calculation is
\[
 \begin{array}{ll}
 \text{node outside }D:&
  \Phi^1_{\bar x}=\mathcal G_{\bar x}(-1)
                                  \otimes\varepsilon(B),\\[2mm]
 \text{tangency to }D:&
  \operatorname{Gr}^a\Phi^1_{\bar x}
      =\operatorname{Gr}^a\mathcal G_{\bar x}^{\rm loc}
                                  \otimes\varepsilon(B),\\[2mm]
 \text{node of }D:&
  \operatorname{Gr}^a\Phi^1_{\bar x}
      =\operatorname{Gr}^a\mathcal G_{\bar x}^{\rm loc}
                                  \otimes\varepsilon(B).
 \end{array}
 \tag{WDB14}
\]
Here the latter two expressions mean the stalks of the constant
extensions of the local graded sheaves of 3.1.2; they do not mean
that the original sheaf was defined at a deleted boundary point.
For the first case the Picard--Lefschetz vanishing cycle supplies
the twist \((-1)\). For tangency, apply vanishing cycles to
\(0\to j_!\overline{\mathbb Q}_\ell\to
\overline{\mathbb Q}_\ell\to
\overline{\mathbb Q}_{\ell D}\to0\): the constant ambient
term has no vanishing cycles, which shifts the divisor computation
by one. At a node of \(D\), the corresponding two-step exact
sequence through the normalization of \(D\) shifts the computation
by two and gives the same degree-one sign representation. These are
3.1.3.1--3.1.5.4.

For each original boundary point of \(X-U\), (WDB5) applied to
\(\mathcal F_0^{[0]}\) gives integral local weights. Hypothesis
(A) makes its graded sheaves locally constant across the puncture.
Taking their external tensor product gives the local filtration of
\(\mathcal G_0\) used in (WDB14); its weights are sums of
integers. At an interior point the weight is zero. The sign
representation has Frobenius eigenvalue \(1\) or \(-1\) and
weight zero, and \((-1)\) adds two to the weight. Every case of
(WDB14) therefore proves
\[
 \text{all Frobenius eigenvalues of }\Phi^1_{\bar x}
 \text{ have integral weights.}
 \tag{WDB15}
\]
The proof retains the distinction between the three geometric
supports while obtaining the same integrality conclusion from them.

Let \(\bar\eta\) be the local geometric generic point at
\(\bar t\). The vanishing of \(\Phi^a\) for \(a\ne1\)
gives the exact sequence 3.2.9.1
\[
 0\to\mathcal R^1_{\bar t}
   \to\mathcal R^1_{\bar\eta}
   \to\Phi^1_{\bar x}
   \to\mathcal R^2_{\bar t}
   \to\mathcal R^2_{\bar\eta}\to0.
 \tag{WDB16}
\]
In particular \(\mathcal R^1\) has no sections supported at
\(\bar t\). Varying \(\bar t\) gives the injective sheaf map
\(\mathcal R^1_0\hookrightarrow w_*w^*\mathcal R^1_0\)
of 3.2.8.

Choose the finite pure-constituent filtration \(G\) on
\(w^*\mathcal R^1_0\) obtained above and extend it by intersection:
\[
 G^a\mathcal R^1_0
 =\mathcal R^1_0\cap w_*G^a(w^*\mathcal R^1_0),
 \qquad \mathcal K_{a,0}=\operatorname{Gr}_G^a\mathcal R^1_0.
 \tag{WDB17}
\]
The intersection is in \(w_*w^*\mathcal R^1_0\).
Write \(b_a\) for the pure weight of \(w^*\mathcal K_{a,0}\).
At \(\bar t\), this is the filtration induced by intersection
with the subspace in (WDB16); consequently taking its associated
graded sequence gives
\[
 0\to (\mathcal K_{a,0})_{\bar t}
 \to (\mathcal K_{a,0})_{\bar\eta}
 \to A_{a,\bar t}\to0,
 \qquad A_{a,\bar t}
      \text{ a subquotient of }\Phi^1_{\bar x}.
 \tag{WDB18}
\]
For completeness, this subquotient statement is elementary for an
inclusion \(E\hookrightarrow V\) and a filtration of \(V\):
filter \(E\) by \(E\cap G^aV\) and \(V/E\) by the images
of \(G^aV\). The associated graded sequence is exact, so the
cokernel in degree \(a\) is a graded piece of \(V/E\).
Here \(V/E\) injects into \(\Phi^1_{\bar x}\) by (WDB16).

If \(A_{a,\bar t}\ne0\), select any one of its Frobenius
eigenvalues \(\lambda\). As a subquotient of the vanishing cycle
it has an integral weight by (WDB15). As a quotient of the local
Weil stalk of the pure sheaf \(w^*\mathcal K_{a,0}\), its weight
lies in \(b_a+\mathbb Z\) by (WDB5). Thus
\[
 A_{a,\bar t}\ne0\quad\Longrightarrow\quad
 w_{N(t)}(\lambda)\in\mathbb Z\cap(b_a+\mathbb Z)
 \quad\Longrightarrow\quad b_a\in\mathbb Z.
 \tag{WDB19}
\]
If instead every \(A_{a,\bar t}\) vanishes, (WDB18) identifies
the special stalk with the entire generic stalk at every boundary
point. The image of specialization is inertia invariant, so inertia
acts trivially and the sheaf is lisse across every such point. It is
therefore lisse on \(P\simeq\mathbb P^1\); its geometric monodromy
is trivial because \(\mathbb P^1\) over an algebraically closed
field is simply connected. It follows that \(\mathcal K_a\) is
geometrically constant. Frobenius on its constant fibre can still
be nontrivial. We have proved the exact alternative of 3.2.9:
\[
 \mathcal K_a\text{ is geometrically constant on }P,
 \quad\text{or}\quad b_a\in\mathbb Z.
 \tag{WDB20}
\]

The pure weight \(b_a\) can be measured at any closed point of
\(W_0\). There its eigenvalues are among those of (WDB13) with
\(b=1\). Applying the strict bound (WDB4) to the weight-zero
coefficient sheaf on that smooth fibre proves \(b_a<2\), which
is 3.2.10. Combining this with (WDB20) gives
\[
 \mathcal K_a\text{ not geometrically constant}
 \quad\Longrightarrow\quad b_a\in\mathbb Z,\quad b_a<2,
 \quad\Longrightarrow\quad b_a\le1.
 \tag{WDB21}
\]
This is the decisive return to the earlier exclusion. A weak bound
\(b_a\le2\) would leave the integer \(2\) available and would
not start the improvement from \(k=0\).

\subsection{Leray transfers the surviving error only once}

If \(\mathcal K_a\) is geometrically constant with fibre \(L_a\),
then
\[
 H^1(P,\mathcal K_a)
   =H^1(\mathbb P^1,\overline{\mathbb Q}_\ell)\otimes L_a=0.
 \tag{WDB22}
\]
The equality is Frobenius-equivariant and holds for every Frobenius
action on \(L_a\). No numerical weight bound for \(L_a\) is
required. In the other case, (WDB8), \(\mathsf B_k\), and (WDB21)
give
\[
 \alpha\in\operatorname{Spec}(F;H^1(P,\mathcal K_a))
 \quad\Longrightarrow\quad
 w_q(\alpha)\le b_a+1+2^{-k}\le2+2^{-k}.
 \tag{WDB23}
\]
The long exact sequences of the finite filtration (WDB17) and the
block-triangular argument following (WDB3) now give
\[
 w_q(\alpha)\le2+2^{-k}
 \quad\text{on }E_2^{11}=H^1(P,\mathcal R^1).
 \tag{WDB24}
\]

For the two outer terms, 3.2.12 uses the elementary degree-zero and
degree-two cohomology of curves from 1.4, not the theorem being
proved. On any fibre, \(H_c^0\) is a space of sections of the
weight-zero lisse sheaf and its eigenvalues have weight zero.
For \(H_c^2\), remove the finitely many singular points and pass
to smooth connected components. Formula 1.4.1.3 identifies the
top cohomology with geometric coinvariants of a stalk, followed by
the twist \((-1)\); its eigenvalues have weight two. If components
are permuted, a finite extension makes them individually defined,
and the weight identity (WDB1) descends the same bound. Hence the
stalk weights of \(\mathcal R^0_0\) and \(\mathcal R^2_0\)
are respectively at most \(0\) and \(2\).

Applying the same elementary curve statements on the base proves
\[
 w_q(\alpha)\le2
 \quad\text{on }E_2^{02}=H^0(P,\mathcal R^2)
 \text{ and }E_2^{20}=H^2(P,\mathcal R^0).
 \tag{WDB25}
\]
For the first assertion, a nonzero global eigensection is nonzero
at some closed-point stalk; Frobenius there acts by the corresponding
power of its global eigenvalue. For the second, removal of finitely
many points does not change \(H_c^2\); on a dense lisse open,
1.4.1.3 expresses it as coinvariants of a stalk with weights at most
zero, twisted by \((-1)\). This adds exactly two. These arguments
also cover constructible sheaves on the base with singular boundary
stalks. The \(\beta\) in the printed explanation of 3.2.12 is
specialized here to the auxiliary weight zero of \(\mathcal G_0\).

Equations (WDB24)--(WDB25) and the spectral sequence (WDB12) prove
\[
 w_q(\mu)\le2+2^{-k}
 \quad\text{for every eigenvalue on }
 H_c^2(\widetilde V,\pi^*\mathcal G).
 \tag{WDB26}
\]
Only one error \(2^{-k}\) has appeared: it enters through the
base cohomology in (WDB23). The fibre coefficient weight has already
been reduced to at most one by the separate integer and strict-bound
argument. Applying \(\mathsf B_k\) twice without that reduction
would instead give \(2+2\cdot2^{-k}\) and no improvement after
taking a square root.

\subsection{The nonzero class survives the two exact injections}

K\"unneth gives a Frobenius-equivariant direct-sum decomposition in
degree two. Its \((1,1)\) summand provides the first map below;
the second is pullback along the blow-up:
\[
 \begin{split}
 H_c^1(U,\mathcal F^{[0]})\otimes
 H_c^1(U,\mathcal F^{[0]})
 &\lhook\joinrel\longrightarrow H_c^2(V,\mathcal G)\\
 &\lhook\joinrel\longrightarrow
 H_c^2(\widetilde V,\pi^*\mathcal G).
 \end{split}
 \tag{WDB27}
\]
The second injection is 3.1.1.4. More precisely 3.1.1.3 identifies
\[
 H_c^2(\widetilde V,\pi^*\mathcal G)
 \simeq H_c^2(V,\mathcal G)
       \oplus H^0(V\cap A,\mathcal G)(-1),
 \tag{WDB28}
\]
with pullback as the inclusion of the first summand. The center is
disjoint from \(D\) by (B), so these are the actual blow-up centers
in the open surface. Equivalently the Poincar\'e-dual transpose
described in 3.1.1.4 supplies a retraction. The extra summand in
(WDB28) is retained rather than identified with zero.

If \(v\ne0\) is an eigenvector of \(F\) with eigenvalue
\(\alpha^{[0]}\) in the original degree-one space, choose a
linear functional \(\lambda\) with \(\lambda(v)=1\). Then
\((\lambda\otimes\lambda)(v\otimes v)=1\), so
\(v\otimes v\ne0\). Both maps in (WDB27) are injective, hence
their composite carries it to a nonzero eigenvector with eigenvalue
\((\alpha^{[0]})^2\). No exterior square is being taken:
the product is on two copies of the curve, so the odd cohomological
degree does not force this tensor to vanish. Invertibility and the
existence of an eigenvector hold over \(\overline{\mathbb Q}_\ell\)
without any diagonalizability assumption.

Applying (WDB26) to this specific image gives
\[
 \begin{split}
 2\bigl(w_q(\alpha)-\beta\bigr)
 &=w_q\bigl((c\alpha)^2\bigr)\le2+2^{-k},\\
 w_q(\alpha)&\le\beta+1+2^{-(k+1)}.
 \end{split}
 \tag{WDB29}
\]
This proves \(\mathsf B_{k+1}\) under (A)--(C), with the original
eigenvalue and original coefficient weight restored exactly.

\subsection{Removing the reduction hypotheses preserves the offending eigenvalue}

For a finite extension \(\mathbb F_{q^d}/\mathbb F_q\), the
geometric cohomology space is unchanged and \(F\) is replaced by
\(F^d\). Thus an eigenvalue \(\alpha\) is replaced by
\(\alpha^d\), and (WDB1) proves equality of its old and new
weights. A bound over the extension gives the identical bound over
the original field. Deligne uses this operation in 3.2.14 to define
over a finite field an embedding and a general-position pencil that
exist over the algebraic closure, and then to make their finitely
many exceptional points rational. This removes (B) and (C). It does
not itself divide any exponent error by \(d\).

To remove (A), the monodromy theorem supplies a finite surjective
cover \(g:U'_0\to U_0\), with \(U'_0\) smooth, for which
\(g^*\mathcal F_0^{[0]}\) has tame unipotent boundary monodromy.
Purity is preserved by pullback: if \(x'\) lies over \(x\) with
residue degree \(e\), its Frobenius eigenvalues are the \(e\)-th
powers of those at \(x\), and \(N(x')=N(x)^e\).
The exact map used to retain cohomology is
\[
 g^*:H_c^a(U,\mathcal F^{[0]})
        \lhook\joinrel\longrightarrow
        H_c^a(U',g^*\mathcal F^{[0]}).
 \tag{WDB30}
\]
Here is a trace construction which includes the multiplicities at
ramification points. On a connected smooth target curve the finite
surjective map is flat of constant degree \(d_g>0\): its finite
local modules are torsion-free over the discrete valuation rings
of the target and hence free. At a geometric point \(u\), the
trace of the constant coefficient sheaf \(E=\overline{\mathbb Q}_\ell\)
is
\[
 \operatorname{Tr}_{g,u}:\bigoplus_{u'\mid u}E\longrightarrow E,
 \qquad(a_{u'})\longmapsto
 \sum_{u'\mid u}
    \operatorname{length}(\mathcal O_{g^{-1}(u),u'})a_{u'}.
 \tag{WDB30a}
\]
The sum of these lengths is \(d_g\). Compatibility with restriction
follows by decomposing the inverse image of an \'etale neighborhood
into open and closed pieces: each piece is finite flat and its
rank is locally constant. Thus (WDB30a) is a sheaf morphism
\(g_*E\to E\), with the unit followed by trace equal to
\(d_g\,\mathrm{id}\). The projection formula tensors this
identity with \(\mathcal F^{[0]}\). Finite maps are proper and
their geometric fibres have no higher \'etale cohomology, so
\(Rg_*=g_*=g_!\). Taking compact-support cohomology gives the
map (WDB30) with retraction \(d_g^{-1}\operatorname{Tr}_g\).
The coefficient field has characteristic zero, so the inverse
\(d_g^{-1}\) exists even if \(\ell\) divides \(d_g\).
The formula includes inseparable lengths as well. Because these
maps are Frobenius-equivariant, they retain the same eigenvalue.
Connected components can be treated after a finite base extension
and descended by (WDB1). The universal \(\mathsf B_{k+1}\) follows.

\subsection{A finite contradiction, followed by the exact dual lower bound}

Induction now proves (WDB2) for every integer \(k\ge0\). Suppose
an actual eigenvalue had
\(w_q(\alpha)=\beta+1+\varepsilon\) with \(\varepsilon>0\).
The specific finite integer
\[
 K=\max\{0,\lfloor\log_2(1/\varepsilon)\rfloor+1\}
 \quad\text{satisfies}\quad 2^{-K}<\varepsilon.
 \tag{WDB31}
\]
Indeed if \(\varepsilon>1\), \(K=0\) works; otherwise
\(K>\log_2(1/\varepsilon)\). At this finite stage (WDB2)
would say
\(\beta+1+\varepsilon\le\beta+1+2^{-K}\), a contradiction.
Consequently
\[
 w_q(\alpha)\le\beta+1
 \quad\text{on }H_c^1(U,\mathcal F).
 \tag{WDB32}
\]
This is the dyadic \(2^{-k}\) bootstrap in Weil II, not the
\(1/k\) error obtained from arbitrarily large products in the
earlier proof of Weil I. Both retain the offending eigenvalue, but
their geometric estimates and induction parameters differ.

Let now \(j:U_0\hookrightarrow X_0\) with \(X_0\) smooth
projective. Equation (WDB8) transfers (WDB32) to
\(H^1(X,j_*\mathcal F)\). Deligne's duality in 3.2.15 is the
perfect Frobenius-equivariant pairing with values in the trivial
coefficient field
\[
 H^1(X,j_*\mathcal F)\ \otimes\
 H^1(X,j_*\mathcal F^\vee(1))
 \longrightarrow\overline{\mathbb Q}_\ell.
\tag{WDB33}
\]
Its perfectness on these actual image groups can also be seen
directly from Poincar\'e duality on \(U\). The compact-to-ordinary
maps for \(\mathcal F\) and \(\mathcal F^\vee(1)\) are
transposes, up to the degree-one cup-product sign. The kernel of
either transpose is the annihilator of the other image. Quotienting
the compactly supported groups by these kernels therefore induces
a nondegenerate pairing on the two images. Their identification
with the two \(j_*\) cohomology groups gives (WDB33), including
its scalar target.
The dual coefficient sheaf has weight \(-\beta-2\). For an
eigenvalue \(\alpha\) of the first factor, \(\alpha^{-1}\)
occurs in the second, whether or not Frobenius is semisimple.
Applying the already proved upper bound to that factor gives
\[
 -w_q(\alpha)=w_q(\alpha^{-1})
 \le(-\beta-2)+1=-\beta-1.
 \tag{WDB34}
\]
Together with the upper bound this proves
\(w_q(\alpha)=\beta+1\). In degrees zero and two, the invariant
and coinvariant descriptions of 1.4, with the Tate twist \((-1)\)
in degree two, give weights \(\beta\) and \(\beta+2\).
Thus all three degrees give 3.2.3:
\[
 w_q(\alpha)=\beta+i
 \quad\text{on }H^i(X,j_*\mathcal F).
 \tag{WDB35}
\]
For degree one this group is also the image of
\(H_c^1(U,\mathcal F)\to H^1(U,\mathcal F)\), as stated at
the start of 3.2.3. Purity is being asserted for this actual image;
the mixed boundary part of all of \(H_c^1(U,\mathcal F)\) has
not been declared pure.

\subsection{The mixed-sheaf theorem retains and controls the boundary}

Theorem 3.3.1 concerns a morphism \(f:X\to Y\) of schemes of
finite type over \(\mathbb Z\), with \(\ell\) invertible as
in Deligne's conventions. It asserts that a mixed constructible
sheaf \(\mathcal F\) of weights at most \(n\) has
\(R^if_!\mathcal F\) mixed of weights at most \(n+i\).
Deligne proves it by reduction to the curve theorem together with
the local boundary calculation. The reductions are by exact
sequences in the sheaf; decomposition into an open set and its
closed complement; base change on the target; and the
Frobenius-equivariant Leray spectral sequence for a composition.
For the sheaf reduction involving a kernel, the preceding term in
the long exact sequence has bound \(n+i-1\), which is at most
\(n+i\); thus that reduction also preserves the claimed bound.
Relative dimension zero contributes only degree zero and preserves
the pointwise weight. Universal homeomorphisms preserve the
\(\acute etale\) theory. Stratification and composition reduce
to relative curves and then to a lisse pure coefficient sheaf.
These are the explicit reductions 3.3.1(a)--(f).

On the generic relative curve one removes finitely many singular
points, takes a finite covering with tame ramification at infinity,
and uses the trace splitting of its coefficient sheaf. A purely
inseparable field extension is allowed in making the smooth
compactification. Deligne then spreads this finite geometric data
over a neighborhood of the generic point of the reduced target.
Reductions 3.3.1(c),(e) make this neighborhood the whole reduced
target and allow the smooth reduced target used in Deligne's
relative boundary calculation; the target reduction handles its
complement. The resulting
relative geometry is
\[
 \begin{gathered}
 X\xrightarrow{j}\overline X\xleftarrow{i}D,
 \qquad f=\bar f j,\\
 \bar f:\overline X\longrightarrow Y
 \text{ smooth projective of relative dimension }1,\\
 D\longrightarrow Y\text{ finite \'etale}.
 \end{gathered}
 \tag{WDB36}
\]
Here \(\mathcal F\) is lisse pure of weight \(n\) on \(X\)
and tamely ramified along \(D\). There is an actual short exact
sequence on the compactification
\[
 0\to j_!\mathcal F\to j_*\mathcal F
   \to i_*B\to0,
 \qquad B=i^*j_*\mathcal F.
 \tag{WDB37}
\]
The boundary stalk is inertia invariants, not the entire generic
stalk. By 1.8.6--1.8.8 its induced monodromy filtration satisfies
\[
 \operatorname{Gr}_a^M B=0\ (a>0),
 \qquad \operatorname{Gr}_a^M B
     \text{ pure of weight }n+a\ (a\le0).
 \tag{WDB38}
\]
The signs here follow from \(V^I\subset\ker N\subset M_0V\);
when a finite residual inertia group remains, its invariant functor
is exact in characteristic zero and retains the same upper indices.
These sheaves and their filtrations commute with passage to the
fibres in (WDB36).

Applying 3.2.3 on each fibre, and for every coefficient embedding
\(\iota\), proves that \(R^a\bar f_*j_*\mathcal F\) is
pointwise pure of weight \(n+a\). Because \(D\to Y\) is finite,
its higher direct images vanish and its direct image carries the
finite filtration (WDB38) to weights at most \(n\). The long exact
sequence of (WDB37) contains, in degrees zero and one,
\[
 \begin{split}
 0&\to f_!\mathcal F\to\bar f_*j_*\mathcal F
   \to(\bar f i)_*B\to R^1f_!\mathcal F\\
  &\to R^1\bar f_*j_*\mathcal F\to0,
 \end{split}
 \tag{WDB39}
\]
and in degrees \(a\ge2\) gives
\(R^af_!\mathcal F\simeq R^a\bar f_*j_*\mathcal F\).
Thus degree zero has weights at most \(n\); degree one is an
extension of the pure sheaf \(R^1\bar f_*j_*\mathcal F\), of
weight \(n+1\), by the boundary quotient
\(\operatorname{coker}(\bar f_*j_*\mathcal F\to(\bar f i)_*B)\),
of weights at most \(n\); and degree two has weight \(n+2\).
This proves the desired upper bound in the reduced
relative-curve case, and the exact reductions give 3.3.1. Boundary
contributions have been included through (WDB37)--(WDB39), with
their weights and maps specified.

For a finite-type \(\mathbb F_q\)-scheme this gives 3.3.4:
\(H_c^i(X,\mathcal F)\) has weights at most \(n+i\) when
\(\mathcal F_0\) is mixed of weights at most \(n\). If in
addition \(X_0\) is smooth of pure dimension \(d\) and
\(\mathcal F_0\) is lisse and mixed of weights at least \(n\),
then \(\mathcal F_0^\vee\) has weights at most \(-n\).
Apply the upper bound to the exact Poincar\'e dual:
\[
 H^i(X,\mathcal F)^\vee
 \simeq H_c^{2d-i}(X,\mathcal F^\vee)(d),
 \qquad
 -n+(2d-i)-2d=-n-i.
 \tag{WDB40}
\]
Duality reverses weights, giving the lower bound \(n+i\) for
\(H^i(X,\mathcal F)\), which is 3.3.5. For a lisse pure sheaf
of weight \(n\), the natural map from compactly supported to
ordinary cohomology is Frobenius-equivariant, and its image is both
a quotient of a space of weights at most \(n+i\) and a subspace
of one of weights at least \(n+i\). The characteristic-polynomial
argument after (WDB3) then proves the exact conclusion 3.3.6:
\[
 \operatorname{im}\bigl(H_c^i(X,\mathcal F)
                 \longrightarrow H^i(X,\mathcal F)\bigr)
 \quad\text{is pure of weight }n+i.
 \tag{WDB41}
\]
The original compact-support and ordinary cohomology groups are
preserved in this statement; purity is proved for their specified
common image.

\subsection{The exact contradiction mechanism available for comparison}

The decisive numerical chain in this proof is
\[
 \begin{gathered}
 \text{nonzero local boundary quotient}
 \ \xRightarrow{\text{(WDB14)--(WDB19)}}\ b_a\in\mathbb Z,\\
 \{b_a\in\mathbb Z,\ b_a<2\}
 \ \Longrightarrow\ b_a\le1,\\
 \text{all local boundary quotients zero}
 \ \xRightarrow{\text{(WDB18)--(WDB22)}}\ H^1(P,\mathcal K_a)=0,\\
 \mathsf B_k\ \xRightarrow{\text{(WDB23)--(WDB26)}}\
   w_q(\mu)\le2+2^{-k},\\
 \alpha^{[0]}\ \xmapsto{\text{(WDB27)}}\ (\alpha^{[0]})^2
 \ \Longrightarrow\ \mathsf B_{k+1}.
 \end{gathered}
 \tag{WDB42}
\]
The word ``boundary'' in this chain names the actual
specialization-cokernel map (WDB18) into vanishing cycles, and the
constant alternative has the actual cohomology-vanishing map
(WDB22). The improvement is proved only after both alternatives
have been treated. The finite integer (WDB31) then excludes every
strict upper violation, and (WDB33)--(WDB34) exclude every strict
lower violation of the pure image weight. Equations
(WDB37)--(WDB41) extend this calculation to mixed coefficients
without deleting their boundary-supported pieces.
